\documentclass{article}
\usepackage{amsmath,amssymb,amsthm}
\usepackage{algorithm}
\usepackage{algpseudocode}
\usepackage{hyperref}
\usepackage{cleveref}
\usepackage{enumitem}
\usepackage{geometry}
\geometry{margin=1in}
\newtheorem{theorem}{Theorem}
\newtheorem{lemma}{Lemma}
\newtheorem{assumption}{Assumption}
\newtheorem{remark}{Remark}
\begin{document}

\section*{Gradient Descent–Ascent (GDA) for Convex–Concave Saddle‑Point Problems}

\section*{Mathematical Formulation}
Let $L:\mathbb{R}^{n}\times\mathbb{R}^{m}\rightarrow\mathbb{R}$ be a \emph{convex–concave} function, i.e.
\[
L(\cdot ,y)\ \text{is convex for every }y,\qquad 
L(x,\cdot )\ \text{is concave for every }x .
\]
We consider the saddle‑point problem
\[
\min_{y\in\mathcal{Y}}\;\max_{x\in\mathcal{X}} L(x,y),
\]
where $\mathcal{X}\subset\mathbb{R}^{n}$ and $\mathcal{Y}\subset\mathbb{R}^{m}$ are closed convex sets.  
Define the gradient operator
\[
\nabla L(x,y)=\begin{pmatrix}
\nabla_{x}L(x,y)\\[2mm]
-\nabla_{y}L(x,y)
\end{pmatrix}.
\]
A \emph{gradient descent–ascent} (GDA) iteration with a constant step size $\eta>0$ is
\begin{align}
x_{t+1}&= \Pi_{\mathcal{X}}\!\bigl(x_{t}+ \eta\,\nabla_{x}L(x_{t},y_{t})\bigr),\label{eq:update-x}\\
y_{t+1}&= \Pi_{\mathcal{Y}}\!\bigl(y_{t}- \eta\,\nabla_{y}L(x_{t},y_{t})\bigr),\label{eq:update-y}
\end{align}
where $\Pi_{\mathcal{S}}$ denotes Euclidean projection onto set $\mathcal{S}$.
The pair $(x_{t},y_{t})$ is the \emph{primal–dual} iterate at time $t$.

\section*{Pseudocode}
\begin{algorithm}[H]
\caption{Gradient Descent–Ascent (GDA)}\label{alg:GDA}
\begin{algorithmic}[1]
\Require Initial point $(x_{0},y_{0})\in\mathcal{X}\times\mathcal{Y}$, step size $\eta>0$, number of iterations $T$.
\For{$t=0,\dots,T-1$}
    \State Compute gradients $g_{x}^{t}= \nabla_{x}L(x_{t},y_{t})$, $g_{y}^{t}= \nabla_{y}L(x_{t},y_{t})$.
    \State $x_{t+1}\gets \Pi_{\mathcal{X}}\!\bigl(x_{t}+ \eta\, g_{x}^{t}\bigr)$.
    \State $y_{t+1}\gets \Pi_{\mathcal{Y}}\!\bigl(y_{t}- \eta\, g_{y}^{t}\bigr)$.
\EndFor
\State \textbf{Output:} $(x_{T},y_{T})$ (or the averaged iterate $\bar{x}_{T}=\frac1T\sum_{t=0}^{T-1}x_{t}$, $\bar{y}_{T}=\frac1T\sum_{t=0}^{T-1}y_{t}$).
\end{algorithmic}
\end{algorithm}

\section*{Dry Run Example}
Consider the simple scalar function $L(x,y)=\tfrac12 (x-y)^{2}$ defined on $\mathcal{X}=\mathcal{Y}=\mathbb{R}$.  
Here $\nabla_{x}L(x,y)=x-y$, $\nabla_{y}L(x,y)=-(x-y)$.  
We start from $x_{0}=0$, $y_{0}=0$, and use step size $\eta=0.1$.

\begin{center}
\begin{tabular}{c|c|c|c|c}
$t$ & $x_{t}$ & $y_{t}$ & $\nabla_{x}L(x_{t},y_{t})$ & $L(x_{t},y_{t})$\\\hline
0 & $0$ & $0$ & $0-0=0$ & $0$\\
1 & $0+0.1\cdot0=0$ & $0-0.1\cdot0=0$ & $0-0=0$ & $0$\\
2 & $0+0.1\cdot0=0$ & $0-0.1\cdot0=0$ & $0$ & $0$\\
3 & $0$ & $0$ & $0$ & $0$
\end{tabular}
\end{center}
Because the initial point already satisfies $x_{0}=y_{0}$, the gradient is zero and the iterates remain at the saddle point.  
If we instead start from $x_{0}=1$, $y_{0}=0$:

\begin{center}
\begin{tabular}{c|c|c|c|c}
$t$ & $x_{t}$ & $y_{t}$ & $\nabla_{x}L$ & $L$ \\ \hline
0 & $1$ & $0$ & $1-0=1$ & $0.5$\\
1 & $1+0.1\cdot1=1.1$ & $0-0.1\cdot(-1)=0.1$ & $1.1-0.1=1.0$ & $0.5\,(1.0)^{2}=0.5$\\
2 & $1.1+0.1\cdot1.0=1.2$ & $0.1-0.1\cdot(-1.0)=0.2$ & $1.2-0.2=1.0$ & $0.5$\\
3 & $1.3$ & $0.3$ & $1.0$ & $0.5$
\end{tabular}
\end{center}
The primal–dual gap $L(x_{t},y_{t})$ stays constant because the function is bilinear; averaging the iterates yields the exact saddle point $(\bar x,\bar y)=(1.2,0.2)$ with gap $0.5$, illustrating the need for averaging to obtain convergence rates (see Section~\ref{sec:theory}).

\section*{Assumptions}\label{sec:assumptions}
\begin{assumption}[Convex–concave structure]\label{as:convexconcave}
$L(\cdot ,y)$ is convex for all $y\in\mathcal{Y}$ and $L(x,\cdot )$ is concave for all $x\in\mathcal{X}$.
\end{assumption}
\begin{assumption}[Smoothness]\label{as:smooth}
$L$ is $L$‑smooth jointly, i.e.
\[
\|\nabla L(u)-\nabla L(v)\|_{2}\leq L\|u-v\|_{2},\qquad 
\forall\,u,v\in\mathcal{X}\times\mathcal{Y}.
\]
Equivalently,
\[
\|\nabla_{x}L(x_{1},y_{1})-\nabla_{x}L(x_{2},y_{2})\|_{2}\le L\|(x_{1},y_{1})-(x_{2},y_{2})\|_{2},
\]
and the same bound holds for $\nabla_{y}L$.
\end{assumption}
\begin{assumption}[Bounded feasible sets]\label{as:bounded}
There exist $D_{x},D_{y}>0$ such that
\[
\|x-x^{\star}\|_{2}\le D_{x},\qquad \|y-y^{\star}\|_{2}\le D_{y},
\]
for any saddle point $(x^{\star},y^{\star})\in\mathcal{X}\times\mathcal{Y}$.
\end{assumption}
\begin{assumption}[Existence of a saddle point]\label{as:saddle}
There exists at least one pair $(x^{\star},y^{\star})\in\mathcal{X}\times\mathcal{Y}$ satisfying
\[
L(x^{\star},y)\le L(x^{\star},y^{\star})\le L(x,y^{\star}),\qquad\forall (x,y)\in\mathcal{X}\times\mathcal{Y}.
\]
\end{assumption}

\section*{Main Convergence Theorem}\label{sec:theory}
\begin{theorem}[Primal–dual gap of averaged GDA]\label{thm:convergence}
Assume \ref{as:convexconcave}–\ref{as:saddle} and choose a constant step size
\[
0<\eta\le \frac{1}{L}.
\]
Define the averaged iterates
\[
\bar{x}_{T}\;:=\;\frac{1}{T}\sum_{t=0}^{T-1}x_{t},\qquad 
\bar{y}_{T}\;:=\;\frac{1}{T}\sum_{t=0}^{T-1}y_{t}.
\]
Then the \emph{primal–dual gap}
\[
\mathcal{G}_{T}\;:=\;\max_{y\in\mathcal{Y}}L(\bar{x}_{T},y)-\min_{x\in\mathcal{X}}L(x,\bar{y}_{T})
\]
obeys
\begin{equation}
\mathcal{G}_{T}\;\le\;\frac{D^{2}}{2\eta T},\qquad 
D^{2}:=D_{x}^{2}+D_{y}^{2}. \label{eq:rate}
\end{equation}
Consequently $\mathcal{G}_{T}=O\!\bigl(\tfrac{1}{T}\bigr)$.
\end{theorem}
\begin{remark}
The rate \eqref{eq:rate} matches the optimal $O(1/T)$ rate for deterministic first‑order methods on smooth convex–concave problems.
\end{remark}

\section*{Theoretical Properties}
\begin{enumerate}[label=\arabic*)]
\item \textbf{Stability.} The condition $\eta\le 1/L$ guarantees that the operator $\mathcal{T}(z)=z-\eta \nabla L(z)$ is \emph{non‑expansive} on the product space, which yields the descent inequality used in the proof.
\item \textbf{Hyper‑parameter sensitivity.} If $\eta$ is chosen larger than $1/L$, the iterates may diverge (the operator becomes expansive). Conversely, a much smaller $\eta$ slows convergence linearly in $1/\eta$ according to \eqref{eq:rate}.
\item \textbf{Comparison to baselines.}
    \begin{itemize}
        \item Standard Gradient Descent on $\min_{y}L(x^{\star},y)$ attains $O(1/T)$ under convexity alone.
        \item Mirror‑prox (extragradient) methods can also achieve $O(1/T)$ without the restriction $\eta\le1/L$, but require an extra gradient evaluation per iteration.
        \item GDA is thus the simplest method attaining the optimal rate when the smoothness constant is known.
    \end{itemize}
\end{enumerate}

\section*{Preliminaries (Definitions and Lemmas)}
Let $z:=\begin{pmatrix}x\\ y\end{pmatrix}\in\mathbb{R}^{n+m}$ and define the monotone operator
\[
F(z):=\begin{pmatrix}
-\nabla_{x}L(x,y)\\[1mm]
\nabla_{y}L(x,y)
\end{pmatrix}
 = -\,\nabla L(x,y).
\]
Because of convex–concave structure, $F$ is \emph{monotone}:
\[
\langle F(z)-F(z'),\,z-z'\rangle \ge 0,\qquad \forall z,z'. \tag{M}
\]
\begin{lemma}[Co‑coercivity of a $L$‑smooth monotone operator]\label{lem:coco}
If $F$ is $L$‑Lipschitz (Assumption~\ref{as:smooth}) and monotone, then for any $z,z'$,
\[
\langle F(z)-F(z'),\,z-z'\rangle \;\ge\; \frac{1}{L}\,\|F(z)-F(z')\|^{2}. \tag{C}
\]
\end{lemma}
\begin{proof}
Standard proof: by $L$‑Lipschitzness,
$\|F(z)-F(z')\|^{2}\le L\,\langle F(z)-F(z'),\,z-z'\rangle$, and rearranging yields \eqref{C}. \qedhere
\end{proof}

\section*{Main Proof (Step‑by‑Step Derivation)}\label{sec:proof}
Denote $z_{t}:=(x_{t},y_{t})$, $z^{\star}:=(x^{\star},y^{\star})$, and the projection operator $\Pi:=\Pi_{\mathcal{X}\times\mathcal{Y}}$ (product of the two Euclidean projections).  
The GDA update \eqref{eq:update-x}–\eqref{eq:update-y} can be compactly written as
\[
z_{t+1}= \Pi\!\bigl(z_{t} - \eta\,F(z_{t})\bigr). \tag{U}
\]

\begin{enumerate}[label=\textbf{[Step \arabic*]}]
\item \textbf{Descent Lemma for the distance to a saddle point.}  
Because projection onto a convex set is non‑expansive,
\[
\|z_{t+1}-z^{\star}\|^{2}
\le \|z_{t}-\eta F(z_{t})-z^{\star}\|^{2}. \tag{1}
\]

\item \textbf{Expand the RHS of (1).}
\begin{align*}
\|z_{t}-\eta F(z_{t})-z^{\star}\|^{2}
&= \|z_{t}-z^{\star}\|^{2}
-2\eta\langle F(z_{t}),\,z_{t}-z^{\star}\rangle
+\eta^{2}\|F(z_{t})\|^{2}. \tag{2}
\end{align*}

\item \textbf{Use monotonicity (M).}  
Since $F(z^{\star})=0$ (first‑order optimality of a saddle point),
\[
\langle F(z_{t})-F(z^{\star}),\,z_{t}-z^{\star}\rangle
= \langle F(z_{t}),\,z_{t}-z^{\star}\rangle \ge 0. \tag{3}
\]

\item \textbf{Apply co‑coercivity (Lemma~\ref{lem:coco}).}
\[
\langle F(z_{t}),\,z_{t}-z^{\star}\rangle
\ge \frac{1}{L}\|F(z_{t})\|^{2}. \tag{4}
\]

\item \textbf{Insert (4) into (2).}
\[
\|z_{t}-\eta F(z_{t})-z^{\star}\|^{2}
\le \|z_{t}-z^{\star}\|^{2}
-2\eta\frac{1}{L}\|F(z_{t})\|^{2}
+\eta^{2}\|F(z_{t})\|^{2}. \tag{5}
\]

\item \textbf{Collect the $\|F(z_{t})\|^{2}$ terms.}
\[
-2\eta\frac{1}{L}+\eta^{2}
= -\eta\Bigl(\frac{2}{L}-\eta\Bigr). \tag{6}
\]

\item \textbf{Choose step size $\eta\le 1/L$.}  
Then $\frac{2}{L}-\eta\ge \frac{1}{L}$, so from (5)–(6)
\[
\|z_{t+1}-z^{\star}\|^{2}
\le \|z_{t}-z^{\star}\|^{2}
-\frac{\eta}{L}\|F(z_{t})\|^{2}. \tag{7}
\]

\item \textbf{Telescoping sum.}  
Sum \eqref{7} for $t=0,\dots,T-1$:
\[
\|z_{T}-z^{\star}\|^{2}
\le \|z_{0}-z^{\star}\|^{2}
-\frac{\eta}{L}\sum_{t=0}^{T-1}\|F(z_{t})\|^{2}. \tag{8}
\]

\item \textbf{Upper bound the left‑hand side by zero.}  
Since the norm is non‑negative,
\[
\frac{\eta}{L}\sum_{t=0}^{T-1}\|F(z_{t})\|^{2}
\le \|z_{0}-z^{\star}\|^{2}\le D^{2},\qquad 
D^{2}=D_{x}^{2}+D_{y}^{2}. \tag{9}
\]

\item \textbf{Relate $\|F(z_{t})\|$ to the primal–dual gap.}  
For any $z=(x,y)$,
\[
\underbrace{L(x,y^{\star})-L(x^{\star},y)}_{\displaystyle =\mathcal{G}(x,y)}
\le \langle F(z),\,z-z^{\star}\rangle
\le \|F(z)\|\,\|z-z^{\star}\|\le \|F(z)\|\,D. \tag{10}
\]
Dividing by $D$ and averaging over $t$ gives
\[
\frac{1}{T}\sum_{t=0}^{T-1}\mathcal{G}(x_{t},y_{t})
\le \frac{D}{T}\sum_{t=0}^{T-1}\|F(z_{t})\|
\le \frac{D}{T}\sqrt{T\sum_{t=0}^{T-1}\|F(z_{t})\|^{2}}
\stackrel{(9)}{\le}\frac{D^{2}}{2\eta T}. \tag{11}
\]

\item \textbf{Jensen’s inequality for the averaged iterates.}  
Because $L(\cdot ,y)$ is convex and $L(x,\cdot )$ is concave,
\[
\max_{y}L(\bar{x}_{T},y)-\min_{x}L(x,\bar{y}_{T})
\le \frac{1}{T}\sum_{t=0}^{T-1}\mathcal{G}(x_{t},y_{t}). \tag{12}
\]

\item \textbf{Combine (11) and (12).}  
Thus the primal–dual gap of the averaged iterates satisfies
\[
\mathcal{G}_{T}
\le \frac{D^{2}}{2\eta T}. \tag{13}
\]

\item \textbf{Conclusion.}  
Inequality \eqref{13} is exactly the bound \eqref{eq:rate} claimed in Theorem~\ref{thm:convergence}. ∎
\end{enumerate}

\section*{Rate Derivation (With Constants)}\label{sec:rate}
The bound \eqref{eq:rate} can be rewritten as
\[
\mathcal{G}_{T}\le \underbrace{\frac{D^{2}}{2}}_{\displaystyle C}\,\frac{1}{\eta T},
\]
where the constant $C=D^{2}/2$ depends only on the diameter of the feasible sets.  
Choosing the largest admissible step size $\eta=1/L$ yields the clean expression
\[
\mathcal{G}_{T}\le \frac{L D^{2}}{2T}. \tag{R}
\]

\section*{Special Cases}
\begin{enumerate}[label=\alph*)]
\item \textbf{Strongly convex–strongly concave case.}  
If $L(\cdot ,y)$ is $\mu_{x}$‑strongly convex and $L(x,\cdot )$ is $\mu_{y}$‑strongly concave, the monotone operator $F$ becomes \emph{strongly monotone}. The same proof with the stronger inequality
\[
\langle F(z)-F(z^{\star}),\,z-z^{\star}\rangle\ge \mu\|z-z^{\star}\|^{2},
\quad \mu:=\min\{\mu_{x},\mu_{y}\},
\]
leads to a linear convergence rate
\[
\|z_{t}-z^{\star}\|^{2}\le (1-\eta\mu)^{t}\|z_{0}-z^{\star}\|^{2}.
\]

\item \textbf{Unconstrained domains.}  
If $\mathcal{X}=\mathbb{R}^{n}$ and $\mathcal{Y}=\mathbb{R}^{m}$, the projection $\Pi$ disappears. The proof simplifies: step~[Step 1] becomes an equality, and the same $O(1/T)$ bound holds under the same step‑size restriction.

\item \textbf{Stochastic gradients.}  
When only unbiased stochastic estimates $\tilde g_{x}^{t},\tilde g_{y}^{t}$ are available, an additional variance term $\sigma^{2}$ appears in (9). The resulting bound is
\[
\mathcal{G}_{T}\le \frac{D^{2}}{2\eta T}+ \frac{\eta\sigma^{2}}{2},
\]
optimally balanced at $\eta\asymp 1/\sqrt{T}$, yielding the classic $O(1/\sqrt{T})$ stochastic rate.

\end{enumerate}

\section*{Discussion}
The presented analysis shows that plain Gradient Descent–Ascent, despite its simplicity, attains the optimal deterministic $O(1/T)$ rate for smooth convex–concave saddle‑point problems when a constant step size not exceeding $1/L$ is used and the output is the \emph{average} of the iterates. The proof hinges on:
\begin{itemize}
\item monotonicity of the saddle‑point operator,
\item $L$‑smoothness giving co‑coercivity,
\item non‑expansiveness of Euclidean projection.
\end{itemize}
These ingredients together produce a clean telescoping inequality (Step~[Step 7]) that directly yields the rate. The result aligns with the known lower bound for first‑order methods and thus establishes the theoretical optimality of GDA in this setting.

\end{document}